\documentclass[12pt, letterpaper]{article}
\usepackage[utf8]{inputenc}
\usepackage[spanish]{babel}
\usepackage{amsmath}
\usepackage{amsfonts}
\usepackage{geometry}
\geometry{top=2cm, bottom=2cm, left=2.5cm, right=2.5cm}
\usepackage{enumitem}

\begin{document}

% Membrete
\noindent
\textbf{Saint Louis School} \\
\textbf{Asignatura:} Matemática \\
\textbf{Nivel:} 7° Básico

\vspace{0.5cm}

\begin{center}
    \Large \textbf{Guía de Práctica - Sesión 1} \\
    \large \textbf{Números Enteros y Potencias}
\end{center}

\vspace{0.5cm}

\noindent
\textbf{Nombre:} \rule{8cm}{0.4pt} \hfill \textbf{Fecha:} \rule{3cm}{0.4pt}

\vspace{0.8cm}

\section*{I. Números Enteros en la Recta Numérica}
\noindent Ubica los siguientes conjuntos de números en una recta numérica imaginaria y luego ordénalos de \textbf{menor a mayor}.

\begin{enumerate}[label=\alph*)]
    \item $\{-5, 2, 0, -8, 4, -1\}$
    \vspace{0.5cm}
    \item $\{15, -12, -2, 7, -20\}$
\end{enumerate}

\vspace{0.5cm}

\section*{II. Adición y Sustracción de Enteros}
\noindent Resuelve las siguientes operaciones aplicando la regla de los signos. Recuerda escribir tu desarrollo si es necesario.

\begin{enumerate}[label=\alph*)]
    \item $(-18) + 25 = $ \vspace{0.4cm}
    \item $14 - (-9) = $ \vspace{0.4cm}
    \item $(-7) + (-12) = $ \vspace{0.4cm}
    \item $-30 - 15 = $ \vspace{0.4cm}
    \item $42 + (-50) = $
\end{enumerate}

\vspace{0.5cm}

\section*{III. Resolviendo Problemas con Enteros}
\noindent Lee atentamente cada situación, plantea la operación matemática correspondiente y responde.

\begin{enumerate}[label=\arabic*.]
    \item Un ascensor se encuentra en el piso 4. Si desciende 7 pisos, ¿en qué piso se encuentra ahora?
    \vspace{1.5cm}
    \item El saldo de una cuenta bancaria es de $\$-15.000$. Si se realiza un depósito de $\$25.000$, ¿cuál es el nuevo saldo de la cuenta?
    \vspace{1.5cm}
    \item Durante una noche de invierno, la temperatura mínima fue de $-6^{\circ}\text{C}$ y la máxima alcanzó los $8^{\circ}\text{C}$. ¿De cuántos grados fue la variación de temperatura durante ese día?
    \vspace{1.5cm}
\end{enumerate}

\section*{IV. Potencias de Base y Exponente Natural}
\noindent Expresa cada multiplicación iterada como una potencia y calcula su valor final.

\begin{enumerate}[label=\alph*)]
    \item $3 \times 3 \times 3 \times 3 = $ \vspace{0.5cm}
    \item $10 \times 10 \times 10 = $ \vspace{0.5cm}
    \item $2 \times 2 \times 2 \times 2 \times 2 \times 2 = $
\end{enumerate}

\vspace{0.5cm}

\noindent Calcula el resultado de las siguientes expresiones que combinan sumas y potencias:

\begin{enumerate}[label=\alph*), start=4]
    \item $4^2 + 5^2 = $ \vspace{1cm}
    \item $10^3 - 2^4 = $ \vspace{1cm}
\end{enumerate}

\vspace{1cm}
\hrule
\vspace{0.5cm}

\begin{center}
    \textbf{Ticket de Salida}
\end{center}

\noindent \textit{Resuelve este breve desafío antes de terminar la clase:} \\

\noindent Si a la temperatura de $-3^{\circ}\text{C}$ le sumamos el resultado de $2^3$, ¿cuál es el número final que obtenemos? Muestra tu cálculo.

\vspace{2cm}

\end{document}